\documentclass[10pt,a4paper]{amsart}
\usepackage[margin=19mm]{geometry}
\usepackage{amsmath,amssymb,mathtools}
\usepackage[scr=rsfs,cal=euler]{mathalfa}
\usepackage{xfrac}
\usepackage[unicode,hidelinks,bookmarks=false]{hyperref}
\newcommand{\ie}{i.e.\ }
\newcommand{\cf}{cf.\ }
\newcommand{\resp}{respectively\ }
\newcommand{\Subsection}{\subsection}
\providecommand{\nobreakdash}{\nobreak\hbox{-}}
\setlength{\emergencystretch}{2em}
\multlinegap0pt
\usepackage{amssymb}
\usepackage{algorithm}\usepackage{algpseudocode}
\usepackage{xcolor}
\newtheorem{proposition}{Proposition}
\title{Coordinate-wise splitting algorithms for ODE simulation via Koopman-Lie product formulas}
\author{Arun Banjara, Ibrahem AlJabea, Theodore Papamarkou,, Frank Neubrander}
\date{Source: arXiv:2506.17524v3 (CC BY 4.0)}
\begin{document}
\maketitle

\noindent\emph{Source and licence.} Coordinate-wise splitting algorithms for ODE simulation via Koopman-Lie product formulas. Version arXiv:2506.17524v3 (CC BY 4.0), distributed by arXiv under the Creative Commons Attribution 4.0 International licence, http://creativecommons.org/licenses/by/4.0/, as recorded in the arXiv licence metadata for this exact version and shown on the abstract page https://arxiv.org/abs/2506.17524v3. Full licence notice: \texttt{licenses/arxiv-2506.17524-CC-BY-4.0.txt}.\par\smallskip

\noindent\textit{Editorial scope.} Counted unit: the article's proposition prop:coordinate-generators (source0.tex lines 401--423). The article's complete original proof of it is delivered with it (source0.tex lines 425--448, one \texttt{proof} environment of 24 lines). No mathematics is written, repaired or re-derived: the statement and the proof are the article's own lines, in the article's own notation, and no retained line is substituted or re-flowed.  No further source-local context is needed: every cross-reference of every retained line resolves inside the two delivered spans.  Nothing was omitted from any retained span, so the package carries no omission record.  No retained line contains a citation, so the wrapper carries no bibliography and no bibliography entry.  No retained line contains a figure environment, an image-inclusion command or drawing code, so no image is delivered and the package declares no asset dependency.  Editorial formatting only, in addition: an AMS \texttt{amsart} preamble that restates the article's own declarations and macros used by the retained lines, namely the environments \texttt{proposition} on separate counters, together with the standard packages those lines need.  The retained environments print in this document as Proposition 1 in order of appearance, whereas the article prints the same items with the numbering of its own full document.  The complete native source of the article as distributed in the arXiv e-print for this exact version is bundled unchanged as \texttt{sources/180241/source0.tex}.
\begin{proposition}
\label{prop:coordinate-generators}
Assume that $F$ generates a unique global, jointly continuous, $\Omega$-invariant flow, and that each frozen vector field $F^{[i]}$ generates a unique global, jointly continuous, $\Omega$-invariant frozen-coordinate flow. Let $g\in C^1(O,\mathbb{C})$, and write $g$ for its restriction to $\Omega$. Suppose that
$
g|_\Omega\in C_b(\Omega),
~
F_i\,\frac{\partial g}{\partial u_i}\Big|_\Omega\in C_b(\Omega)$ for $i=1,\ldots,N$, and $F\cdot\nabla g\big|_\Omega\in C_b(\Omega)$. Then $g$ belongs to the domains of $\mathcal{K}$ and $\mathcal{K}_i$, and
\[
\mathcal{K}_i g(u)
=
F_i(u)\frac{\partial g}{\partial u_i}(u),
~
i=1,\ldots,N.
\]
Moreover,
\[
\mathcal{K}g(u)
=
F(u)\cdot\nabla g(u)
=
\sum_{i=1}^N \mathcal{K}_i g(u).
\]
\end{proposition}

\begin{proof}
We give the argument for $\mathcal{K}_i$; the proof for $\mathcal{K}$ is the same, with $\Phi_i$ replaced by $\varphi$. For $u\in\Omega$, the chain rule gives
\[
\frac{d}{ds}g(\Phi_i(s,u))
=
F_i(\Phi_i(s,u))
\frac{\partial g}{\partial u_i}(\Phi_i(s,u)).
\]
Hence,
\[
\frac{T_i(t)g(u)-g(u)}{t}
=
\frac{1}{t}\int_0^t
F_i(\Phi_i(s,u))
\frac{\partial g}{\partial u_i}(\Phi_i(s,u))\,ds .
\]
Since $\Phi_i$ is jointly continuous, the integrand converges uniformly on compact subsets of $\Omega$ to $F_i(u)\frac{\partial g}{\partial u_i}(u)$ as $t\to 0^+$. The assumed boundedness gives the bounded-difference-quotient condition in the definition of the bi-continuous generator. Therefore,
\[
\mathcal{K}_i g(u)
=
F_i(u)\frac{\partial g}{\partial u_i}(u).
\]
Summing these identities gives the stated formula for $\mathcal{K}g$.
\end{proof}

\end{document}
